% !TeX root = ../main.tex

\chapter{基于多智能体图构建与图采样的推荐方法}
\label{cha:multi-agent-graph}

第三章将外部知识以语义向量的形式注入推荐，本章则从图结构的角度注入知识：由多个智能体协同完成知识连接与语义连接的构建，并对图结构进行动态采样优化，最终通过多视角图卷积生成用户与物品表征。

% TODO: 本章方法基于已申请发明专利“基于多智能体图构建和图采样的推荐方法和系统”。
%       按学术规范，在本章首页加脚注说明该章方法的专利关联。

\section{引言}
\label{sec:graph-intro}

如第~\ref{cha:related-work}~章所述，知识图谱增强推荐的关键瓶颈有二：其一是图的构建高度依赖人工，难以覆盖开放域知识，且构建完成后难以随业务演化而更新；其二是图上的长尾节点交互稀疏、头部节点噪声冗余，导致表征学习不充分或过平滑。

本章以“多智能体协同”应对第一个瓶颈：知识融合智能体负责从外部知识库与预训练语言模型中抽取实体关系，语义增强智能体负责在向量空间中建立语义相近节点之间的连接，二者共同将原本单一的交互图扩展为交互、知识、语义三个视角的多视角图。针对第二个瓶颈，由采样优化智能体对长尾节点补全伪链路以缓解稀疏，对头部节点筛选强链路以抑制噪声。在此基础上执行多视角图卷积，并引入分散正则化缓解表征同质化。

\section{问题定义与总体框架}
\label{sec:graph-problem}

给定用户集合 $\mathcal{U}$、物品集合 $\mathcal{I}$ 与交互集合 $\mathcal{O}\subseteq\mathcal{U}\times\mathcal{I}$，本方法将推荐任务建模为多视角图上的节点表征学习问题。总体框架包含四个阶段：多视角图构建、图采样优化、多视角图卷积、双模式推荐。

% TODO: 插入本章方法总体框架图（图 4-1），并在正文中逐段解释四个阶段的数据流向。

\section{多智能体协同的多视角图构建}
\label{sec:multi-view-graph}

多视角图由三类连接构成：交互连接 $\mathcal{E}_{int}$、知识连接 $\mathcal{E}_{kn}$ 与语义连接 $\mathcal{E}_{sem}$。所有连接以元组形式存储，支持增量更新，从而避免全量重建。

\subsection{知识融合智能体}
\label{subsec:knowledge-agent}

知识融合智能体调用外部知识库与预训练语言模型，为用户与物品生成属性标签与语义描述，并识别用户—知识实体、物品—知识实体之间的关联关系，形成知识连接：

\begin{equation}
  e_{kn} = \left\langle \mathit{id}_1,\ r,\ \mathit{id}_2,\ \mathit{conf} \right\rangle ,
  \label{eq:knowledge-tuple}
\end{equation}

其中 $r$ 为关系类型，$\mathit{conf}$ 为置信度，由来源可靠性与模型输出的一致性共同决定。

\subsection{语义增强智能体}
\label{subsec:semantic-agent}

语义增强智能体将用户与物品的语义描述经文本编码模型转换为向量并存入向量数据库，随后检索相似度超过阈值的候选集合，经推理评估后建立三类语义连接（用户—用户、物品—物品、用户—物品）：

\begin{equation}
  e_{sem} = \left\langle \mathit{id}_1,\ \mathit{id}_2,\ \mathit{sim},\ \mathit{type} \right\rangle .
  \label{eq:semantic-tuple}
\end{equation}

与知识连接相比，语义连接不要求共享显式实体，因而能够覆盖开放域中的隐含相似关系。

\subsection{多视角图与增量更新}
\label{subsec:incremental}

三类连接合并构成多视角图 $\mathcal{G} = \{\mathcal{G}_{int}, \mathcal{G}_{kn}, \mathcal{G}_{sem}\}$。当新物品入库或用户产生新交互时，仅需触发相关智能体对局部连接进行增删，无需重建整图。

\section{面向长尾分布的图采样优化}
\label{sec:graph-sampling}

\subsection{长尾节点伪链路补全}
\label{subsec:pseudo-link}

对交互次数低于阈值的用户与物品，采样优化智能体借助向量数据库按语义相似度召回候选节点，并生成伪链路：

\begin{equation}
  e_{pse} = \left\langle u,\ i,\ \mathit{pseudo} \right\rangle ,
  \label{eq:pseudo-tuple}
\end{equation}

伪链路仅参与表征学习，不参与最终的监督标签计算，以避免引入标签噪声。

\subsection{头部节点强链路筛选}
\label{subsec:strong-link}

对交互率高于阈值的头部用户与物品，采样优化智能体依据交互质量与相关性推理，过滤冗余与噪声交互，仅保留强链路：

\begin{equation}
  e_{keep} = \left\langle u,\ i,\ \mathit{strong} \right\rangle .
  \label{eq:keep-tuple}
\end{equation}

该策略在降低图规模的同时抑制了热门物品带来的过平滑问题。

\section{多视角图卷积与表征融合}
\label{sec:multi-view-gcn}

\subsection{视角子图卷积}
\label{subsec:view-gcn}

在每个视角子图上执行图卷积，以交互视角为例：

\begin{equation}
  \mathbf{e}_u^{(l+1)} = \mathrm{AGG}\left(\left\{\mathbf{e}_v^{(l)} \mid v\in\mathcal{N}_u^{int}\right\}\right),
  \label{eq:gcn}
\end{equation}

其中 $\mathrm{AGG}(\cdot)$ 可采用均值聚合或注意力聚合，$\mathcal{N}_u^{int}$ 为用户在交互视角下的邻域。知识视角与语义视角采用相同形式。

\subsection{多视角表征融合}
\label{subsec:view-fusion}

各视角输出按可学习权重融合：

\begin{equation}
  \mathbf{e}_u = \sum_{m\in\{int,kn,sem\}} w_m\,\mathbf{e}_u^{(m)},
  \qquad \sum_{m} w_m = 1 .
  \label{eq:view-fusion}
\end{equation}

\subsection{分散正则化}
\label{subsec:dispersion}

多视角聚合易导致节点表征趋同。为此在训练目标中引入分散正则化项，惩罚表征两两之间过高的余弦相似度，以保持节点表征的差异性。

% TODO: 补充分散正则化的具体形式与权重，并与专利权利要求 6 对应。

\section{双模式推荐}
\label{sec:dual-mode}

\subsection{直接模式}
\label{subsec:direct-mode}

直接模式下，以融合后的用户表征与物品表征的内积作为推荐分数，适用于低时延、高吞吐的在线场景。

\subsection{智能体模式}
\label{subsec:agent-mode}

智能模式在直接模式的基础上引入两个额外环节：兴趣调研智能体根据近期上下文生成动态用户兴趣向量，与静态融合表征融合；排序智能体对初筛列表进行二次排序。该模式精度更高但时延更大，适用于对效果敏感的场景（如首页首屏）。

\section{实验}
\label{sec:graph-experiments}

\subsection{数据集与预处理}
\label{subsec:graph-datasets}

% TODO: 采用 Amazon-Book、LastFM、Yelp2018 等公开数据集，按 KGAT/KGIN 的划分方式构造训练/验证/测试集，
%       并统计用户数、物品数、交互数、知识三元组数与稀疏度。

\subsection{对比方法}
\label{subsec:graph-baselines}

选取 NGCF、LightGCN、KGAT、KGIN、SGL 作为对比方法，覆盖图协同过滤、知识图谱增强与图采样三类技术路线。

\subsection{评价指标与实现细节}
\label{subsec:graph-metrics}

采用 Recall@20 与 NDCG@20 评价 Top-K 推荐效果，并报告训练时间与推理时延。

\subsection{整体性能对比}
\label{subsec:graph-overall}

% TODO: 插入主结果表（三个数据集 × Recall@20/NDCG@20），并分析提升来源。

\subsection{消融实验}
\label{subsec:graph-ablation}

% TODO: 插入消融表，逐项移除：知识视角、语义视角、伪链路补全、强链路筛选、分散正则化，
%       与专利权利要求一一对应，验证各模块必要性。

\subsection{长尾与效率分析}
\label{subsec:graph-longtail}

% TODO: 按交互次数对用户分组（如 <10、10–50、>50），报告各组的 Recall 提升；
%       对比直接模式与智能模式的时延与效果差异。

\section{系统落地与运行效果}
\label{sec:graph-deployment}

% TODO: 本节对应横向项目的系统实现，注意脱敏。
% 建议结构：
%   (1) 部署形态：多视角图的离线构建与增量更新流程；
%   (2) 模式切换：直接模式与智能模式的灰度策略；
%   (3) 运行效果：以相对指标表述，并说明统计口径与观测周期。

\section{本章小结}
\label{sec:graph-summary}

本章提出了一种基于多智能体图构建与图采样的推荐方法。该方法以多智能体协同实现多视角图的自动化与增量式构建，以采样优化缓解长尾稀疏与头部噪声，以多视角图卷积与分散正则化提升表征质量，并提供直接模式与智能体模式两种部署形态。实验与工程实践表明，该方法在推荐精度、长尾覆盖与系统可维护性方面均具有优势。
